\documentclass{beamer}

\usetheme{Madrid}

%================= FONT & LANGUAGE =================
\usepackage[utf8]{inputenc}
\usepackage[T5]{fontenc}
\usepackage[vietnamese]{babel}

%================= PACKAGE =================
\usepackage{amsmath}
\usepackage{tikz}
\usepackage{listings}
\usepackage{xcolor}
\usepackage{tcolorbox}
\usepackage{fontawesome5}

%================= COLOR =================
\definecolor{main1}{RGB}{102,126,234}
\definecolor{main2}{RGB}{118,75,162}
\definecolor{boxbg}{RGB}{255,255,255}

\setbeamercolor{frametitle}{fg=white,bg=main1}
\setbeamercolor{title}{fg=white,bg=main2}

%================= BACKGROUND =================
\setbeamertemplate{background canvas}{
\begin{tikzpicture}[remember picture,overlay]
\path[left color=main1!20,right color=main2!20]
(current page.south west) rectangle (current page.north east);
\end{tikzpicture}
}

%================= BOX =================
\tcbset{
  colback=boxbg,
  colframe=main1,
  arc=3mm,
  boxrule=0.8pt
}

%================= CODE =================
\lstset{
  language=C++,
  basicstyle=\ttfamily\tiny,
  keywordstyle=\color{main2},
  commentstyle=\color{green!50!black},
  breaklines=true
}

%================= INFO =================
\title{Thuật toán DFS (Depth-First Search)}
\author{Lê Văn Hoàng An}
\date{\today}

\begin{document}

\frame{\titlepage}

%================= INTRO =================
\section{Giới thiệu}

\begin{frame}{DFS là gì?}
\begin{tcolorbox}
\begin{itemize}
\item \faSearch DFS = Depth-First Search (Duyệt theo chiều sâu)
\item \faArrowDown Đi sâu nhất có thể trước khi quay lui
\item \faRedo Backtracking khi không còn đường đi
\item \faCogs Sử dụng đệ quy hoặc stack
\end{itemize}
\end{tcolorbox}

\vspace{0.2cm}
\begin{tcolorbox}
\textbf{Ý nghĩa:}
\begin{itemize}
\item Giống khám phá mê cung
\item Đi một hướng đến cùng rồi quay lại
\end{itemize}
\end{tcolorbox}
\end{frame}

%================= GRAPH =================
\section{Minh họa}

\begin{frame}{Minh họa DFS}
\begin{columns}

\column{0.5\textwidth}
\begin{tikzpicture}[scale=1.1, every node/.style={circle, draw}]
\node (A) at (0,0) {A};
\node (B) at (2,1) {B};
\node (C) at (2,-1) {C};
\node (D) at (4,1) {D};
\node (E) at (4,-1) {E};

\draw (A)--(B);
\draw (A)--(C);
\draw (B)--(D);
\draw (C)--(E);
\end{tikzpicture}

\column{0.5\textwidth}
\begin{tcolorbox}
\begin{itemize}
\item \faPlay Bắt đầu từ A
\item \faArrowRight Đi sâu: A → B → D
\item \faUndo Quay lui → A → C → E
\end{itemize}
\end{tcolorbox}

\end{columns}
\end{frame}

%================= THEORY =================
\section{Lý thuyết}

\begin{frame}{Ý tưởng + Giả mã DFS}
\begin{tcolorbox}
\begin{itemize}
\item \faBullseye Mục tiêu: duyệt sâu theo nhánh
\item \faLayerGroup Cấu trúc: stack / đệ quy
\item \faCheckCircle Mỗi đỉnh chỉ thăm 1 lần
\end{itemize}
\end{tcolorbox}

\vspace{0.2cm}

\begin{tcolorbox}
\textbf{Quy trình:}
\begin{enumerate}
\item Đánh dấu đỉnh hiện tại
\item Xử lý đỉnh
\item Duyệt các đỉnh kề chưa thăm
\item Gọi DFS đệ quy
\end{enumerate}
\end{tcolorbox}
\end{frame}

%================= CODE =================
\section{Cài đặt}

\begin{frame}[fragile]{Code DFS (C++)}
\begin{tcolorbox}
\begin{lstlisting}
#include <bits/stdc++.h>
using namespace std;

vector<vector<int>> adj;
vector<bool> visited;

void dfs(int u) {
    visited[u] = true;

    // Xu ly dinh u
    cout << u << " ";

    for (int v : adj[u]) {
        if (!visited[v]) {
            dfs(v);
        }
    }
}

/*
Phan tich:
- Thoi gian: O(V + E)
- Moi dinh duyet 1 lan
- Su dung de quy don gian
*/
\end{lstlisting}
\end{tcolorbox}
\end{frame}

%================= EXAMPLE =================
\section{Ví dụ}

\begin{frame}{Ví dụ minh họa DFS}
\begin{columns}

\column{0.5\textwidth}
\begin{tikzpicture}[scale=1.1, every node/.style={circle, draw}]
\node (1) at (0,0) {1};
\node (2) at (2,1) {2};
\node (3) at (2,-1) {3};
\node (4) at (4,1) {4};
\node (5) at (4,-1) {5};

\draw (1)--(2);
\draw (1)--(3);
\draw (2)--(4);
\draw (3)--(5);
\end{tikzpicture}

\column{0.5\textwidth}
\begin{tcolorbox}
\begin{itemize}
\item \faRoute Thứ tự:
\[
1 \rightarrow 2 \rightarrow 4 \rightarrow 3 \rightarrow 5
\]
\item \faExchangeAlt Có backtracking
\item \faProjectDiagram Không theo lớp như BFS
\end{itemize}
\end{tcolorbox}

\end{columns}
\end{frame}

%================= ANIMATION =================
\section{Animation DFS}

\begin{frame}{Minh họa DFS từng bước}
\centering

\begin{tikzpicture}[
    scale=1.2,
    every node/.style={circle, draw, minimum size=8mm},
    visited/.style={fill=green!40},
    current/.style={fill=yellow!60}
]

\node (1) at (0,0) {1};
\node (2) at (2,1) {2};
\node (3) at (2,-1) {3};
\node (4) at (4,1) {4};
\node (5) at (4,-1) {5};

\draw (1)--(2);
\draw (1)--(3);
\draw (2)--(4);
\draw (3)--(5);

\only<1>{\node[current] at (1) {1};}
\only<2>{\node[visited] at (1) {1}; \node[current] at (2) {2};}
\only<3>{\node[visited] at (1) {1}; \node[visited] at (2) {2}; \node[current] at (4) {4};}
\only<4>{\node[visited] at (1) {1}; \node[visited] at (2) {2}; \node[visited] at (4) {4}; \node[current] at (3) {3};}
\only<5>{\node[visited] at (1) {1}; \node[visited] at (2) {2}; \node[visited] at (4) {4}; \node[visited] at (3) {3}; \node[current] at (5) {5};}

\end{tikzpicture}

\vspace{0.5cm}
\only<1>{Bắt đầu từ 1}
\only<2>{Đi sâu đến 2}
\only<3>{Tiếp tục đến 4}
\only<4>{Quay lui → sang 3}
\only<5>{Hoàn tất DFS}

\end{frame}

%================= APPLICATION =================
\section{Ứng dụng}

\begin{frame}{Ứng dụng của DFS}
\begin{tcolorbox}
\begin{itemize}
\item \faProjectDiagram Tìm thành phần liên thông
\item \faSync Kiểm tra chu trình
\item \faBrain Backtracking, AI, game
\end{itemize}
\end{tcolorbox}
\end{frame}

%================= END =================
\begin{frame}{Kết luận}
\begin{tcolorbox}
\begin{itemize}
\item DFS là thuật toán nền tảng
\item Độ phức tạp tối ưu $O(V+E)$
\item Quan trọng trong nhiều bài toán nâng cao
\end{itemize}
\end{tcolorbox}

\centering
\vspace{0.5cm}
\Huge{\textbf{Cảm ơn!}}
\end{frame}

\end{document}
